% 29_body.tex — 산출물 ㉙ PR/FAQ · 실험 사전등록서 · 베타 진입 기준 (빈 템플릿 / 완성 예시 공용 본문) — gen_product_templates.py 가 Product_specs.py 에서 생성
% 드라이버: 29_PRFAQ실험사전등록_빈템플릿.tex (\wbblanktrue) / 29_PRFAQ실험사전등록_완성예시.tex (\wbblankfalse — 워크북 §X.5 확정 후)
\documentclass[10pt,a4paper]{article}
\usepackage[a4paper,margin=16mm,top=14mm,bottom=20mm]{geometry}
\def\DocNum{29}\def\DocDay{6}\def\DocIdx{4}
\def\DocTitle{PR/FAQ · 실험 사전등록서 · 베타 진입 기준}
\def\DocSub{Press Release/FAQ, Experiment Pre-registration \& Beta Criteria}
\def\DocVariant{\B{빈 템플릿}{딥게이지 완성 예시}}
\usepackage{wbstyle}
\def\DocCirc{㉙}   % newunicodechar(㉛~㊵ 폴백) 뒤에 정의해야 활성 문자로 읽힌다
\begin{document}
\docheader
 {\Bg{[회사명 · 제품명]}{딥게이지 · GAUGE}}
 {\Bg{[v1.x / D+n]}{v1.0 / 2026-08-07 (D+158)}}
 {\Bg{[작성자 — CPO]}{윤하경 CPO\rd{7mm}}}
 {\Bg{[검토자 — 전원]}{전원\rd{7mm}}}

%% ------------------------------------------------------------------ 1
\secbar{1. 보도자료 (출시일에 낼 글을 지금 쓴다)}
\guide{제목 / 부제 / 요약 / 문제 / 해법 / 리더 인용 / 고객 인용 / 시작 방법. 각 칸 3줄 이내. 기능 목록이 아니라 고객의 변화를 쓴다.}
{\footnotesize
\begin{tabularx}{\linewidth}{|L{30mm}|Y|}\hline
\hd{요소} & \hd{내용} \\ \hline
\ifwbblank
\rd{10mm}제목 &  \\ \hline
\rd{10mm}부제 &  \\ \hline
\rd{10mm}요약(한 문단) &  \\ \hline
\rd{10mm}문제 &  \\ \hline
\rd{10mm}해법 &  \\ \hline
\rd{10mm}리더 인용 &  \\ \hline
\rd{10mm}고객 인용(가상) &  \\ \hline
\rd{10mm}시작 방법 &  \\ \hline
\else
제목 & 안산 세영정공, 검사기록 작성 하루 47분 → 19분 — 딥게이지 'GAUGE' 정식 출시 \\ \hline
부제 & 검사원이 찍은 사진과 한마디 음성에서 검사기록·부적합 패키지·8D 초안이 만들어진다. 자동차 2·3차 협력사 3개사 도입 \\ \hline
요약(한 문단) & 12월 27일 정식 출시. 라인에서 찍은 사진과 음성 메모가 검사기록이 되고, 부적합은 사진·기록·판정 패키지로 1차사 제출 초안과 8D 초안이 된다. 파일럿 3사에서 기록 47 → 19분, 통보 3.7일 → 1영업일 이내 \\ \hline
문제 & 라인이 도는 동안 기록할 수 없어 종이 → 엑셀 → 포털로 세 번 옮겨 적는다(31분/일). 부적합 통보 3.7일, 8D 11.5시간, 클레임 2,700만 \\ \hline
해법 & 검사원의 행동은 바뀌지 않는다 — 이미 찍던 사진, 한마디. 기록은 그 자리에서, 근거(누가·언제·어느 사진·어느 기준)는 로그로. 판정은 검사원이 한다 \\ \hline
리더 인용 & '3년 동안 개선안을 냈고 3년 동안 반려됐습니다. 검사원이 문제라고 느끼지 않는 것을 우리가 문제라고 우기지 않는 제품을 만들려 했습니다.' — 강민수 대표 \\ \hline
고객 인용(가상) & '통보가 늦는 건 자료 모으느라였는데, 이제 부적합 등록하면 사진이랑 기록이 붙어서 나와요. 8D는 여전히 우리가 쓰지만, 붙이는 건 안 해요.' — [세영정공] 정미영 품질팀장 \\ \hline
시작 방법 & 라인당 월 45만 원, 좌석 무제한. 검사항목 마스터(엑셀)를 보내면 10영업일 안에 시작 \\ \hline
\fi
\end{tabularx}}

%% ------------------------------------------------------------------ 2
\secbar{2. 고객 FAQ 12문}
\guide{고객이 실제로 물을 것 — 가격, 데이터, 온프렘, 오판정 시 책임, 기존 MES와의 관계, 검사원 교육.}
{\footnotesize
\begin{tabularx}{\linewidth}{|L{8mm}|Y|Y|}\hline
\hd{\#} & \hd{질문} & \hd{답(3줄)} \\ \hline
\ifwbblank
\rd{8mm} &  &  \\ \hline
\rd{8mm} &  &  \\ \hline
\rd{8mm} &  &  \\ \hline
\rd{8mm} &  &  \\ \hline
\rd{8mm} &  &  \\ \hline
\rd{8mm} &  &  \\ \hline
\rd{8mm} &  &  \\ \hline
\rd{8mm} &  &  \\ \hline
\rd{8mm} &  &  \\ \hline
\rd{8mm} &  &  \\ \hline
\rd{8mm} &  &  \\ \hline
\rd{8mm} &  &  \\ \hline
\else
1 & 얼마입니까? & 라인당 월 45만, 계정 무제한. 구축 300~500만 일회성. 최소 라인 2 · 12개월(㉘ 항목 7) \\ \hline
2 & 우리 도면이 클라우드로 올라갑니까? & 검사 사진·기록만. 도면 이미지·규격 문서는 올리지 않는 '도면 미업로드 모드'(㉘ 항목 6) \\ \hline
3 & 온프렘으로 됩니까? & 지금은 아님. 유료 5개사 달성 또는 1차사 채널 계약 시 재검토(㉖ 항목 5) — '검토 중'이 아니라 이 조건 \\ \hline
4 & AI가 틀리면 누가 책임집니까? & 최종 판정은 검사원. 초안은 참고. 계약상 당사 책임은 연 계약액의 10\% 상한(㉗·㉘ 항목 6) \\ \hline
5 & 자동으로 합격/불합격을 내나요? & 아닙니다. 검사원이 승인해야 판정. 자동판정 기능 없음 \\ \hline
6 & MES와 연동됩니까? & 엑셀·PDF 내보내기로 연결. MES 연동은 범위 밖(㉖ Won't) \\ \hline
7 & 1차사 포털에 바로 올라갑니까? & H사 규격 제출 파일 생성. 업로드는 팀장님이(2분). 다른 규격은 검토 대상 \\ \hline
8 & 검사원 교육은? & 1회 2시간, 세 동작. 전 검사원 출신 담당이 첫 주 동행 \\ \hline
9 & 오후에 사진이 잘 안 읽힌다는데? & 조도 낮으면 앱이 경고·재촬영 안내. 정확도는 데이터셋 기준으로 공개(㉗ 항목 3) \\ \hline
10 & 우리 데이터를 학습에 씁니까? & 비식별 검사 사진·기록만. 도면 제외, 계약서 명시 \\ \hline
11 & 해지는? & 30일 전 통지, 잔여 50\% 환불. 구축비 환불 없음 \\ \hline
12 & 라인 전용 검사항목을 넣어 줄 수 있나요? & 매핑 범위 안에서 가능. 구조 변경 요청은 범위 결정서를 거쳐 서면 답변(㉖ 항목 5 규칙) \\ \hline
\fi
\end{tabularx}}

%% ------------------------------------------------------------------ 3
\secbar{3. 내부 FAQ 15문}
\guide{무엇이 어렵나 / 왜 지금 / 무엇을 안 하나 / 단가와 원가 / 온프렘 요구에 대한 입장 / 실패하면 어떻게 아나. 답하지 못하는 질문은 '미확정 — 실험 X-ID'.}
{\footnotesize
\begin{tabularx}{\linewidth}{|L{8mm}|Y|Y|}\hline
\hd{\#} & \hd{질문} & \hd{답 또는 미확정 근거} \\ \hline
\ifwbblank
\rd{8mm} &  &  \\ \hline
\rd{8mm} &  &  \\ \hline
\rd{8mm} &  &  \\ \hline
\rd{8mm} &  &  \\ \hline
\rd{8mm} &  &  \\ \hline
\rd{8mm} &  &  \\ \hline
\rd{8mm} &  &  \\ \hline
\rd{8mm} &  &  \\ \hline
\rd{8mm} &  &  \\ \hline
\rd{8mm} &  &  \\ \hline
\rd{8mm} &  &  \\ \hline
\rd{8mm} &  &  \\ \hline
\rd{8mm} &  &  \\ \hline
\rd{8mm} &  &  \\ \hline
\rd{8mm} &  &  \\ \hline
\else
1 & 무엇이 가장 어려운가? & 오후 조도의 OCR(84.7\%)과 검사원의 수용(A-05b) — 모델이 아니라 현장의 문제 \\ \hline
2 & 왜 지금인가? & 1차사의 포털 제출 요구 시작(I-08·23), VLM 단가 32원 \\ \hline
3 & 무엇을 안 하는가? & 온프렘·MES·자동판정·포털 2종 이상·항목 구조 커스텀(㉖ 항목 5) \\ \hline
4 & 단가 45만은 원가를 감당하는가? & 표준 18.1만(13.4\%), 헤비 69.5만(51.5\%). 헤비가 몇 곳인지는 미확정 — 계측 없음 \\ \hline
5 & 온프렘 요구가 오면? & Won't 재검토 조건을 답한다. 동보프레스는 도면 미업로드 모드로(D+200) \\ \hline
6 & 실패하면 어떻게 아는가? & X-04·05·06·07 사전등록(항목 4) + 오류 예산 소진(㉗ 항목 5) \\ \hline
7 & 검사원이 초안을 받아들이는가? & 미확정 — X-04(D+290). X-01에서 5명 중 2명 별도 종이 기록 \\ \hline
8 & 공장이 라인 과금을 받아들이는가? & 미확정 — X-05(각 사 유료 전환일) \\ \hline
9 & 조도 경고가 재촬영을 줄이는가? & 미확정 — X-07(D+280). 안 줄면 오후 라인 초안 제외 \\ \hline
10 & 유료 5개사(D+360)는 가능한가? & 미확정 — 관찰만. 3사 + 파이프라인 4사. 실험이 아니라 세일즈 \\ \hline
11 & FN 8.1\%(미합격)로 GA해도 되는가? & 기준 유지, 배포 조건을 낮춘다 — '합격 전 베타', 승인 필수, 세영 오전 1라인(㉗ 항목 3) \\ \hline
12 & 조현우 이후 개발은? & 17 mm 재산정(㉖ v1.1). 남은 5.5로 Must 1.1·Should 1.5·예비 2.9 \\ \hline
13 & 구축비를 ARR에 넣는가? & 넣지 않는다(㉘ 항목 6). 4,860 / 6,060 — 보고서에 두 숫자를 다 쓴다 \\ \hline
14 & 12월 청구서 780만은? & (D+284 추가) 초과 570 중 동보 444. 가설 4 — 계측 먼저(㉘ 항목 4) \\ \hline
15 & 이 보도자료에서 아직 참이 아닌 문장은? & '통보 1영업일 이내'(실측 없음) · '3개사 도입'(D+158 유료 0). D+300 전에 참이 되거나 빠진다 \\ \hline
\fi
\end{tabularx}}

%% ------------------------------------------------------------------ 4
\landscapeon
\secbar{4. 실험 사전등록서 (실험 1건 = 1행)}
\guide{OEC(단일 판정 지표) · 가설 · 최소 표본과 산출 근거 · 분석 시점 · SRM 점검 · peeking 금지 선언 · 성공/중단 기준. 결과를 보기 전에 서명한다.}
{\footnotesize
\begin{tabularx}{\linewidth}{|L{14mm}|Y|L{34mm}|L{34mm}|L{26mm}|Y|L{26mm}|}\hline
\hd{X-ID} & \hd{가설(A-ID)} & \hd{OEC} & \hd{최소 표본·근거} & \hd{분석 시점} & \hd{성공 / 중단 기준} & \hd{서명·일자} \\ \hline
\ifwbblank
\rd{13mm} &  &  &  &  &  &  \\ \hline
\rd{13mm} &  &  &  &  &  &  \\ \hline
\rd{13mm} &  &  &  &  &  &  \\ \hline
\rd{13mm} &  &  &  &  &  &  \\ \hline
\else
X-04 판정 초안 베타 & A-05b · A-11 & 초안 수용률(수정 없이 승인) ≥ 70\% 그리고 회피 행동 < 10\% & 판정 600건 = 검사원 3 × 시간대 2 × 100, 설계효과 3 & D+290 고정 · SRM: 검사원별 ± 20\% & 성공: 둘 다 → B-17 재검토 / 불확정: 수용률만 → 초안 유지 / 중단: 회피 20\% 초과 → 초안 OFF & 3인 D+158 \\ \hline
X-05 라인 과금 수용 & A-06 & 계약 라인 ÷ 실사용 ≥ 0.8 그리고 좌석 요구 0건 & 3사 전수(모집단 3) & 각 사 유료 전환일 & 성공 3/3 / 중단: 1사라도 좌석 요구 → ㉘ 항목 1 재검토 & 3인 D+158 \\ \hline
X-06 온보딩 2주 & A-08 & 매핑 완료까지 영업일 ≤ 10 & 3사 전수(217 · 62 · 183개) & 계약 후 15영업일 & 성공 3/3 / 중단: 20일 초과 → B-05 재설계 & 3인 D+158 \\ \hline
X-07 조도 경고 효과 & A-02(부분) & 오후 3시 이후 재촬영률 18\% → 8\% 미만(ON vs OFF) & arm당 177 → 200 사진. 두 비율 검정, α .05, 검정력 .8 & D+280 고정 · SRM 55:45 ± 5 & 성공 < 8\% → 전 라인 ON / 실패 ≥ 12\% → 오후 초안 제외 / 중단: 기록 시간 +5분 & 3인 D+158 \\ \hline
\fi
\end{tabularx}}
\landscapeoff

%% ------------------------------------------------------------------ 5
\secbar{5. 베타 진입 기준 — 알파 / 베타 / GA}
\guide{각 단계의 진입 조건(숫자), 판정자, 되돌림 조건. ㉖ 항목 7과 같은 숫자.}
{\footnotesize
\begin{tabularx}{\linewidth}{|L{20mm}|Y|L{30mm}|Y|}\hline
\hd{단계} & \hd{진입 조건(숫자)} & \hd{판정자} & \hd{되돌림 조건} \\ \hline
\ifwbblank
\rd{11mm}알파 &  &  &  \\ \hline
\rd{11mm}베타 &  &  &  \\ \hline
\rd{11mm}GA &  &  &  \\ \hline
\else
알파 & 내부 + 세영 1라인 · 자동 생성 성공률 ≥ 90\% · 크래시 0건/주 · 매핑 완료 & CTO & 성공률 < 80\% 2주 연속 → 내부로 \\ \hline
베타 & 3사 전 라인 · 기록 ≤ 25분/일 · 오전 OCR ≥ 90\% · 재촬영 < 15\% · 초안은 세영 오전 1라인만 & CPO & 기록 > 35분 또는 재촬영 > 25\% 1사 이상 → 알파로 \\ \hline
GA & 유료 3사 · eval v1 측정 · 오류 예산 가동 · HITL · 자동판정 OFF · 미합격 시 '합격 전 베타' & CEO(전원) & 오류 예산 2개월 연속 소진 또는 2사 해지 통지 → 베타로 \\ \hline
\fi
\end{tabularx}}

\end{document}